\section{An open region of K3 periods}
\label{deg:section}

Theorem~\ref{sew:chain} permits the centers of successive annuli to
move.  We now construct a degeneration for which every possible new
center admits another block.  This is the reason an infinite chain
can be continued after the exactness adjustments.

\begin{proposition}\label{deg:open}
There is a nonempty open subset of the full marked K3 period domain
such that every surface with period in this subset has property $L$
at every point.
\end{proposition}

The openness here is unrestricted: the nearby surfaces need not be
projective.  We shall first produce one smooth quartic with a finite
collection of buffered chart constructions, then deform all those
constructions to its full local deformation space.

\subsection{The semistable model and its regular pencils}

Choose two general smooth quadrics $Q_+,Q_-\subset\mathbb P^3$.
Their intersection $E$ is a smooth elliptic quartic.  For a general
quartic form $G$ consider the family
\begin{equation}\label{deg:quartic-family}
                         Q_+Q_-=\alpha G.
\end{equation}
Here the same symbols $Q_+,Q_-$ denote their quadratic equations.
The sixteen zeros $Z$ of $G|_E$ are simple.  Near such a zero the
total space has equation $xy=\alpha t$, up to holomorphic changes
of coordinates.  Blowing up one component, locally the ideal
$(x,\alpha)$, is a small resolution.  For example, in one chart
write $\alpha=xa$; the strict transform has $y=at$ and the map to
the parameter disc is $\alpha=xa$.  In the other chart write
$x=\alpha b$; the equation becomes $t=by$, with smooth coordinates
$(\alpha,b,y)$ and projection $\alpha$.  Thus the resolved total
space is smooth and the central fiber has normal crossings.

Its two components, with a suitable choice of labels, are a quadric
$S^+$ and the blowup $S^-$ of a quadric at $Z$; they meet along the
strict transform of $E$, again denoted $E$.  The resolution is
crepant because it is small and the original hypersurface is
Gorenstein.  Adjunction for the quartic family supplies a nowhere-zero
relative dualizing form.  On $S^\pm$ it becomes a logarithmic
symplectic form $\omega_\pm$ with simple pole along $E$ and opposite
nonzero residues.  In a normal-crossing chart we have
\begin{equation}\label{deg:neck-form}
       uv=\alpha,\qquad
       \omega_\alpha=k(u,v,t)\,\frac{du}{u}\wedge dt,
\end{equation}
where $k$ is holomorphic and nowhere zero and $t$ is a coordinate
along $E$.  In particular $k(0,0,t)dt$ is its residue on the
$u$-component; the residue on the $v$-component is its negative.
For small nonzero $\alpha$, the fibers of
Equation~\eqref{deg:quartic-family} are smooth quartic K3 surfaces.

Call a fiber of a base-point-free pencil on $S^+$ or $S^-$
\emph{regular} if it is a smooth rational curve, the pencil map is
smooth along it, and it meets $E$ transversely at two distinct
points.  We use the two ruling pencils, pulled back on $S^-$, and
on $S^-$ also the pencils of $(1,1)$ curves through two distinct
points of $Z$, with their two base points subtracted.

\begin{lemma}\label{deg:regular-pencils}
The quadrics and $G$ can be chosen so that:
\begin{enumerate}
 \item Through every point of $E$ on each component there is a
       regular fiber of one of these finitely many pencils.
 \item The union of their regular fibers contains the complement
       of a proper algebraic subset of each component.
 \item Every regular fiber has, away from $E$, a transverse
       intersection with a regular fiber of another pencil.
\end{enumerate}
\end{lemma}

\begin{proof}
On a plain quadric, identify the two rulings with the projections
of $\mathbb P^1\times\mathbb P^1$.  Their restrictions to $E$ have
degree two.  They cannot both ramify at one point: otherwise the
differential of the inclusion of the smooth curve $E$ in the
quadric would vanish.  This proves the first assertion on $S^+$.

For the blown-up quadric, a ruling fiber is additionally forbidden
when it passes through a point of $Z$.  Choose $Z$ generally so
that its points avoid both ruling ramification sets, and so
that a forbidden level of one ruling and a forbidden level of the
other have a common point on $E$ only at a point of $Z$.  To check
that these are proper conditions, write the two projections as
$\pi_1,\pi_2$.  Their ramification sets are disjoint.  Conditions
involving one ramification set and one $\pi_i(Z)$ exclude only
finitely many values for the corresponding point of $Z$.  For
distinct $z_i,z_j$, the condition that the grid point with
coordinates $(\pi_1(z_i),\pi_2(z_j))$ lie on $E$ is a proper
condition on the pair.  For $i=j$ the grid point is $z_i$ itself.
Thus, away from $Z$, at least one ruling remains regular.

Fix $z_i\in Z$.  Choose another center $z_j$ and take the plane
through the tangent line to $E$ at $z_i$ and through $z_j$.  For
general choices its intersection with the quadric is a smooth
$(1,1)$ curve.  Its divisor on $E$ is
\[
                        2z_i+z_j+z_k,
\]
with $z_k$ different from $z_i,z_j$, and it passes through no other
point of $Z$.  Its strict transform is a member of the pencil of
$(1,1)$ curves through $z_i,z_j$.  The chord through these two
points is not a ruling line; hence the pencil has exactly these
two simple base points, which are resolved by the blowups.  An
order-two contact at $z_i$ becomes a simple transverse
intersection with the strict transform of $E$.  The simple
intersection at $z_j$ disappears, and the intersection at $z_k$
remains.  The selected strict transform is therefore a regular
fiber through the point of $E$ over $z_i$.

All the generic requirements just used are simultaneous finite
proper conditions.  More explicitly, on the elliptic curve the
residual $z_k$ is determined by the hyperplane divisor class and
$2z_i+z_j$; excluding coincidences or another center imposes a
proper condition on at most three of the centers.  The other
requirements involve at most two centers and fixed branch data.
The divisor $Z$ ranges over the complete degree-sixteen linear
system $|\mathcal O_E(4)|$.  Indeed the restriction of quartic
forms is surjective, as follows from the Koszul resolution for
the intersection of two quadrics.  In an ordered general divisor
of this degree any specified fewer than sixteen points vary
freely; the remaining points satisfy the one elliptic sum
condition.  Hence the finitely many proper conditions above can
all be avoided.

For each of the finitely many pencils, only finitely many members
fail to be smooth or meet $E$ nontransversely.  Their union is a
proper algebraic set, proving the second assertion.  Finally, on
any regular rational fiber, at least one of the two ruling
projections is nonconstant.  This is immediate for a ruling fiber
by using the other ruling; for a $(1,1)$ fiber both projections are
nonconstant.  Away from finitely many points of the fiber the
corresponding ruling member is regular and the projection has
nonzero derivative.  Such an intersection is transverse and can
be chosen off $E$.
\end{proof}

We now make finite compact choices, each with an open buffer.
Let $G_{\rm reg}$ be the union of \emph{all} regular fibers from
the listed pencils.  Its complement is contained in a proper
algebraic set, and $G_{\rm reg}$ contains a neighborhood of $E$.
By the first assertion of Lemma~\ref{deg:regular-pencils}, compact
subfamilies of regular fibers supply entries at every point of
$E$.  By its third assertion each required family can, after
restriction, be equipped with a transverse partner at an
interior point, with a positive transversality margin in local
charts.  Compactness supplies finitely many such choices near
$E$.

For each point of the algebraic exceptional set away from $E$,
choose a coordinate line not contained in that set and a small
radius whose boundary avoids it.  The boundary then lies in
$G_{\rm reg}$.  Small translations and tilts retain this boundary
containment, so the disc construction supplies a neighborhood of
possible centers.  Points of $G_{\rm reg}$ themselves have regular
coordinate neighborhoods.  On the compact region outside a
chosen neighborhood of $E$, finitely many of these two kinds of
center neighborhoods suffice.  Only after choosing these finitely
many discs do we add finitely many regular patches covering their
compact boundaries, along with the corresponding regular
families and transverse partners.  Thus every retained datum
belongs to a finite buffered collection.  The algebraic
exceptional set is used before the finite selection.  These
disc choices will be used with
Corollary~\ref{loc:disc} after constructing the strips.

\subsection{Core and neck coordinates}

\begin{lemma}\label{deg:core-coordinates}
On a small parameter disc of any compact regular family there are
coordinates $(z,p)\in\mathbb P^1\times\Delta_r$ in which the two
ends are $z=0,\infty$ and
\[
              \omega_\pm=\frac{1}{2\pi i}\frac{dz}{z}\wedge dp.
\]
Either end can be designated as the entrance.  With entrance at
$\infty$, exit at $0$, and $z=c\exp(2\pi i w)$, this reads
$\omega_\pm=dw\wedge dp$.  The parameter $p$ is a local affine
coordinate for either end position on $E$, read using its signed
residue differential.
\end{lemma}

\begin{proof}
Over a small disc the smooth proper family of rational fibers is
a holomorphic $\mathbb P^1$ bundle and is trivial there.  One can
see this by choosing a local section, taking its relative
degree-one line bundle, and using sections and base change to
identify the family with a projective bundle.  The two distinct
end sections can then be put at zero and infinity.

In a parameter $s$ on this disc write the logarithmic form as
$f(z,s)dz\wedge ds$.  For each $s$, the one-form
$f(z,s)dz$ on $\mathbb P^1$ has only simple poles at zero and
infinity.  It is therefore $b(s)dz/z$.  Nondegeneracy makes $b$
nonzero.  Set $dp=2\pi i b(s)ds$.

Let $t_+(s)$ and $t_-(s)$ be the exit and entrance positions on
$E$, and let $\rho$ be the residue differential on this
component.  Taking residues gives the exact identities
\[
             dp=2\pi i\,t_+^*\rho=-2\pi i\,t_-^*\rho.
\]
They also follow from the constancy of the elliptic sum of the
two intersection points.  Since both ends are transverse,
$t_+$ and $t_-$ are locally invertible.  Thus $p$ moves each
end freely and the stated residue interpretation follows.
\end{proof}

In a fixed normal-crossing chart at an end, its nonzero normal
coordinate on the corresponding central component has the
uniform expansions
\begin{align}
 u&=k_-(p)z^{-1}\bigl(1+O(z^{-1})\bigr)
                                    &&\text{at the entrance},
                                      \label{deg:entrance}\\
 u&=k_+(p)z\bigl(1+O(z)\bigr)
                                    &&\text{at the exit}.
                                      \label{deg:exit}
\end{align}
The coordinate along $E$ differs from its end value by
$O(z^{-1})$ or $O(z)$ respectively.  The functions $k_\pm$
are nonzero, with uniformly bounded logarithmic derivatives on
the selected smaller parameter discs.  These estimates, with
any fixed number of derivatives in logarithmic and parameter
coordinates, are uniform over the finite buffered atlas.

Choose a small positive normal cutoff $d$.  Truncate a core at
its two ends where $|u|$ is comparable to $d$, retaining fixed
buffers in logarithmic height.  In the coordinate $w$ of
Lemma~\ref{deg:core-coordinates}, this is a cylinder of length
tending to infinity as $d\to0$, uniformly over the compact
families.  For every fixed $d$ it is a relatively compact annulus
times a disc in the smooth part of the central fiber.
Lemma~\ref{loc:deformation} deforms this map into nearby smooth
fibers on a slightly larger annulus and parameter disc.  Its
symplectic normalization gives maps with pullback form
$dw\wedge dp$ converging to the central maps on these fixed
domains.  This uses only finite-domain deformations, and can be
done simultaneously for the finite raw atlas.  Restrictions and
recentering of its larger parameter discs then supply all the
moving centers used below.

The connecting neck has a more direct description.  In
Equation~\eqref{deg:neck-form} put
\[
 u=c\exp(2\pi i w),\qquad v=\alpha/u,
\]
and allow it to run from $|u|\asymp d$ to $|v|\asymp d$,
with fixed additional logarithmic buffers.  Normalize its
transverse coordinate by
\begin{equation}\label{deg:neck-normalization}
 p(w,t)=\int_{t_0}^t2\pi i\,k(u,\alpha/u,s)\,ds.
\end{equation}
This gives $dw\wedge dp=\omega_\alpha$.  On this entire neck,
including its buffers, both $|u|$ and $|v|$ are at most $Cd$.
Consequently
\[
 p(w,t)=\int_{t_0}^t2\pi i\,k(0,0,s)\,ds+O(d).
\]
The estimate is uniform in neck length, and holds with any fixed
number of derivatives on smaller parameter discs and logarithmic
buffers.  Indeed each $w$ derivative differentiates $u$ or $v$
by a fixed constant multiple, and both remain $O(d)$; parameter
derivatives are bounded on the fixed $t$-patches.  The derivative
in $t$ is bounded away from zero, so the normalization is
uniformly invertible there.

\begin{figure}[ht]
\centering
\begin{tikzpicture}[x=0.93cm,y=0.85cm,>=Latex,font=\small]
 \fill[black!7] (0,0) rectangle (2.8,1.05);
 \fill[black!3] (2.4,0) rectangle (8.6,1.05);
 \fill[black!7] (8.2,0) rectangle (11,1.05);
 \draw (0,0) rectangle (2.8,1.05);
 \draw (2.4,0) rectangle (8.6,1.05);
 \draw (8.2,0) rectangle (11,1.05);
 \draw[dashed] (2.6,-0.2) -- (2.6,1.35);
 \draw[dashed] (8.4,-0.2) -- (8.4,1.35);
 \node at (1.2,0.53) {$S^+$ core};
 \node at (5.5,0.72) {neck: $uv=\alpha$};
 \node at (9.8,0.53) {$S^-$ core};
 \draw[->] (3.5,0.25) -- (7.5,0.25);
 \node[above] at (2.6,1.35) {$|u|=d$};
 \node[above] at (8.4,1.35) {$|v|=d$};
 \node[below] at (5.5,-0.1) {$Y$ increases; $|u|$ decreases and $|v|$ increases};
\end{tikzpicture}
\caption{The core--neck joins in logarithmic height, not to scale.
The overlap widths are fixed.  Once $d$ is fixed, the neck length
grows as $\alpha\to0$, while its transverse width stays bounded
below.  After sewing, the same parameter $p$ labels a leaf across
both joins.}
\label{deg:neck-figure}
\end{figure}

Figure~\ref{deg:neck-figure} separates the fixed-width overlaps from
the part of the cylinder whose length grows.  Only the former
enter the local transition estimates.

\begin{lemma}[Buffered seams]\label{deg:seams}
There are fixed positive parameter and seam widths and a fixed
bound on translation derivatives with the following property.
Choose $d$ sufficiently small, then $\alpha\neq0$ sufficiently
small depending on $d$.  At matching nominal end positions the
normalized core and neck charts have degree-one symplectic
transitions of the form
\[
         F_j=J_j\circ(\operatorname{id}+E_j),\qquad
         J_j(w,p)=(w+a_j(p),p),
\]
on buffered seams.  After choosing long-coordinate origins and
cut heights, they satisfy the domain hypotheses of
Theorem~\ref{sew:chain}, apart from exactness, and
\[
                     \sup_j\|E_j\|\leq C d+\epsilon(d,\alpha),
              \qquad \epsilon(d,\alpha)\longrightarrow0
                            \quad(\alpha\to0, d\text{ fixed}).
\]
The widths and the bound for $a_j'$ can be chosen before $d$.
\end{lemma}

\begin{proof}
At an exit seam Equation~\eqref{deg:exit} has leading part
$u=k_+(p)z$.  Taking a local logarithm of the nonzero coefficient
gives a change $w_{\rm neck}=w_{\rm core}+a(p)$.  Its derivative
is bounded by the logarithmic derivative of $k_+$.  Its transverse
change is exactly the identity at leading order: both parameters
are primitives of the same signed residue differential, with
the same constant at the matching nominal end.

At the entrance into the next component, the leading relation is
$v=k_-(p')/z'$ and $uv=\alpha$.  Hence $z'$ is a nonzero
holomorphic multiple of $u$, so again $w'=w+a(p)$, with the same
orientation.  The sign of the residue changes between the two
components, and the use of $z'=\infty$ at the entrance changes
it once more.  The transverse primitives therefore again agree
as germs.  The potentially large constant $\log\alpha$ affects
the translation itself, not its derivative in $p$.

The actual seam maps differ from these leading maps by $O(d)$
from the expansions and Equation~\eqref{deg:neck-normalization},
plus the core deformation error $\epsilon(d,\alpha)$.  These
estimates are in $(\log u,t)$ or $(\log v,t)$ coordinates.  The
leading changes are uniformly invertible there.  The inverse
function theorem with the reserved buffers therefore produces
the actual inverse branches throughout the seams.  Local branches
agree by uniqueness near the specified leading branch, including
after going once around the period.  Their lifts have degree one,
and they are symplectic because both charts pull back the same
surface form to $dw\wedge dp$.

For completeness, the cuts can be prepared without imposing
incompatible origins at the two ends of a block.  Work first in
raw coordinates.  Enlarge each core range to
$cd<|z|<1/(cd)$ with a fixed sufficiently small $c>0$, and add
fixed logarithmic buffers to every neck.  At a seam choose the
two matching cut heights by its leading formula at the nominal
center.  Each raw block has independent room for the choice of
its two cuts.  Then translate its single long-coordinate origin
so that its entrance is at height zero; its exit defines $L_j$.
For small $d$, and then small $\alpha$, all core and neck lengths
are at least one.  Shrink the fixed transverse discs so that
the variation of $\operatorname{Im}a_j$ about its central value
is less than $h/8$.  The logarithmic derivative bounds show that
this shrinking is independent of $d$ and $\alpha$.  The two
cut heights then give precisely
$|\operatorname{Im}a_j+L_j|<h/8$ with the required margins.
\end{proof}

\subsection{Infinite continuation and transverse partners}

We spell out the compactness statement needed for an infinite
construction.  The finite atlas has larger and smaller end-position
patches whose interiors cover $E$.  At a designated end choose a
regular family and a neck patch whose smaller patches contain it.
Their larger patches give a uniform positive buffer, after taking
a finite refinement if necessary.  Thus a transverse translation
whose size is less than a fixed threshold still belongs to the
chosen larger patches, for every designated position on $E$.
Both orderings of the ends of a core are included in the atlas.
If a later choice requires a smaller positive working radius,
use more centers from these same compact families.  Their smaller
end-position neighborhoods again cover $E$, and a finite subcover
can be chosen.  These are restrictions and translations of the
same raw charts, so the uniform derivative and buffer bounds do
not change.  In particular, the positive parameter width retained
by the sewing theorem can be fixed before the cutoff $d$ without
losing coverage.

Start with any prescribed core or neck chart.  Match its outgoing
nominal end with a new chart.  By Lemma~\ref{deg:seams} its period
defect is $O(d+\epsilon(d,\alpha))$.  Apply
Lemma~\ref{sew:centers} to make the seam exact, translating the
actual and model transverse coordinates of the newly added chart
together.  Use that updated model's outgoing end for the next
choice.  It is some point of $E$, where another uniformly
buffered choice is available.  Continue alternately through
cores on the two components and intervening necks.  In the
opposite direction perform the corresponding adjustment on
each newly added input block.  The initial block need not be
moved in either construction.

The displacement of the nominal centers over many steps can be
large.  This causes no loss of the estimates: at each step the
comparison is with that block's own translated model on a
larger prepared patch, and the outgoing model position is
updated before the following seam is formed.  The set of
possible positions is the whole compact curve $E$, not a
single fixed local parameter disc.  The actual parameter disc
in each centered chart, in contrast, has the same fixed radius.
This is exactly the distinction used in
Lemma~\ref{sew:centers}.

Take $d$ small and then $\alpha$ small enough that all seam
errors, including the period adjustments, are below the threshold
of Theorem~\ref{sew:chain}.  We obtain periodic symplectic
immersions of $\C\times\Delta_b$ whose corrected charts are
uniformly close to the corresponding raw charts.  By making the
initial errors smaller, their first derivatives on the fixed
interior core patches are as close as required.

At a designated interior crossing use an additional chain starting
from the transverse partner pencil.  In the central model the
two tangent directions have a fixed positive angle in the chosen
local coordinates.  The two corrected families still have
independent directions there, and the partner chart still covers
a smaller neighborhood of the crossing.  For every sufficiently
small transverse parameter on the first strip, a point on its
leaf in this core patch therefore lies on a leaf of the partner
family, with transverse intersection.  This uses the partner
\emph{family}: no interpolation at a prescribed point is needed.

By Remark~\ref{loc:alignment} and Proposition~\ref{loc:criterion},
these strips imply property $L$ at every covered point on the
first strip, including points in other blocks of its chain.
The unbounded coordinate may be translated to the crossing,
and the common parameter $p$ identifies the same leaf all along
the chain.  A strip initialized in a neck reaches one of the
adjacent prescribed regular core families and uses its chosen
partner in exactly this fashion.

\subsection{Coverage and unrestricted deformation}

There are two coverage issues: the middle of an increasingly long
neck and the neighborhood of either cutoff.  Neither may be
discarded as a small set.  We verify them in coordinates with
uniform bounds.

Inside a core away from its ends, the deformed chart and its
sewing correction are close in $C^1$ to the central regular
chart.  Local invertibility with a margin implies that its image
contains every prescribed smaller regular patch, under a local
identification of the smooth fibers.  This includes the compact
boundary patches of all the coordinate discs chosen earlier.
The same statement is uniform over compact subfamilies and
small recenterings.

For a neck, the normalized raw map is already a coordinate chart
in $(\log u,t)$ on the whole buffered interval from $|u|\asymp d$
to $|v|\asymp d$.  The correction supplied by
Theorem~\ref{sew:chain} is uniformly close to the identity in
$(w,p)$ and is defined with fixed margins, independently of
the interval length.  Its image contains a uniformly smaller
source cylinder: at any point of that smaller cylinder, a
fixed-radius coordinate ball lies in the larger one, and the
small-change inverse theorem applies there.  Thus the corrected
chart covers the entire smaller neck interval and the smaller
end-position patch, not only a compact part chosen in advance.

At the $u$-cut use a normal-crossing patch about a point of $E$.
On the component side, Equations~\eqref{deg:entrance} and
\eqref{deg:exit} give uniformly nonsingular leading coordinates
$(\log u,t)$, with error $O(|u|)$.  The coordinate $t$ differs
from the endpoint coordinate by $O(|u|)$.  Choose a fixed small
normal size $d_*>0$, smaller than a fixed fraction of the
available end-position width.  The selected core charts cover
the region from a fixed multiple of $d$ up to $d_*$, after
shrinking their end-position patches.  Their sewing corrections
are uniformly small in the same logarithmic coordinates.
The neck chart covers below that multiple of $d$ and overlaps
the core region by its extra logarithmic buffers.  The identical
argument applies at the $v$-cut.  A finite collection of the
smaller end-position patches covers $E$.

This argument does not ask for a lower bound on an ordinary
metric chart radius at $|u|=d$.  All inverse estimates there
are in logarithmic and end-position coordinates, where the
leading maps and the buffers have uniform bounds.  One first
chooses the position widths and the needed error bound for
coverage and for the core crossings, and then chooses $d$.
The fixed $d_*$ can be reduced after the widths are chosen;
finally take $d$ much smaller than $d_*$.  On the remaining
compact region outside these end neighborhoods, the preceding
finite collection of regular patches and coordinate discs
applies.  If the end neighborhoods have been reduced in choosing
$d_*$, the intervening closed annular region is already covered
by the same regular core families; it is a fixed compact region
before $d$ is chosen.  Properness of the semistable family ensures that
these core and neck regions cover every nearby smooth fiber.

At each reserved coordinate disc away from $E$, use the
submersion coordinates of the total space to transport its
local coordinate construction to the smooth fibers.  Its
boundary, and the boundaries of sufficiently nearby translated
or tilted discs, remain in the regular patches where $L$ was
just proved.  Corollary~\ref{loc:disc} gives $L$ throughout the
remaining smaller coordinate neighborhoods.  Thus $L$ holds
at every point of the selected smooth quartic.

Now fix the small nonzero $\alpha$ used above.  Every raw core
and neck chart is a map from a \emph{finite} annulus times a
disc, with margins.  There are finitely many raw charts;
all the moving versions are restrictions and translations of
these larger charts.  Apply Lemma~\ref{loc:deformation} in the
full local deformation family of this smooth K3, choosing a
nearby holomorphic symplectic form and normalizing each chart.
They deform simultaneously with arbitrarily small additional
error.  Their seam inverse branches persist: on each compact
buffered seam they are the nearby branches of the original
local inverse, and uniqueness makes those branches agree
throughout the seam and around its period circle.  Thus all
the domain, derivative, period-adjustment, coverage, and
transversality estimates persist in a sufficiently small
unrestricted deformation neighborhood.  This argument does
not deform the infinitely many corrected charts individually;
it deforms the finite raw atlas and repeats its uniformly
controlled construction.

The order of choices is worth recording because it prevents
any dependence of the allowable error on the length of the
neck:
\begin{enumerate}
 \item Choose the finite compact regular families, partner
       crossings, end-position patches, disc boundaries,
       parameter widths and logarithmic buffers.
 \item Choose an allowable sewing and coverage error for
       these data; the derivative bounds in the leading
       shears have already been fixed.
 \item Choose $d$ so that the $O(d)$ errors are below that
       allowance, and then choose a nonzero $\alpha$ so that
       all fixed-$d$ deformation errors are below it.
 \item Choose the full deformation neighborhood of that
       fixed smooth quartic so that its additional errors
       remain within the same allowance.
\end{enumerate}

K3 deformation theory gives a smooth local deformation space,
and the marked period map is a local isomorphism by local
Torelli; see \cite{Huybrechts}.  Its image of this last
neighborhood is a nonempty open subset of the full period
domain.  Any other marked K3 with the same period has an
isomorphic underlying surface: signs and reflections in
algebraic roots match the K\"ahler chambers while fixing the
period, and global Torelli then applies \cite{Huybrechts}.
Consequently every surface with a period in this open subset
has $L$ everywhere, proving
Proposition~\ref{deg:open}.  Theorem~\ref{ana:cap} consequently
gives convex approximation for each of these surfaces.
